\documentclass[10pt,a4paper]{amsart}
\usepackage{amsmath,amssymb,amsthm,graphicx}
% Retained minimal preamble for the Stacks Project chapter more-morphisms.tex (More on Morphisms).
% Verbatim excerpt of the chapter preamble: Xy-pic setup (native preamble lines 14-18) and the native
% theorem declarations with the mathematical macros (native preamble lines 154-217). Copied mechanically;
% nothing added, removed or redefined. The wrapper class and the amsmath/amssymb/amsthm/graphicx base
% packages are supplied by the assembler, so the class/externaldocument part of the native preamble is omitted.
% For commutative diagrams we use Xy-pic
\usepackage[all]{xy}

% We use 2cell for 2-commutative diagrams.
\xyoption{2cell}
% Theorem environments.
%
\theoremstyle{plain}
\newtheorem{theorem}[subsection]{Theorem}
\newtheorem{proposition}[subsection]{Proposition}
\newtheorem{lemma}[subsection]{Lemma}

\theoremstyle{definition}
\newtheorem{definition}[subsection]{Definition}
\newtheorem{example}[subsection]{Example}
\newtheorem{exercise}[subsection]{Exercise}
\newtheorem{situation}[subsection]{Situation}

\theoremstyle{remark}
\newtheorem{remark}[subsection]{Remark}
\newtheorem{remarks}[subsection]{Remarks}

\numberwithin{equation}{subsection}

% Macros
%
\def\lim{\mathop{\mathrm{lim}}\nolimits}
\def\colim{\mathop{\mathrm{colim}}\nolimits}
\def\Spec{\mathop{\mathrm{Spec}}}
\def\Hom{\mathop{\mathrm{Hom}}\nolimits}
\def\Ext{\mathop{\mathrm{Ext}}\nolimits}
\def\SheafHom{\mathop{\mathcal{H}\!\mathit{om}}\nolimits}
\def\SheafExt{\mathop{\mathcal{E}\!\mathit{xt}}\nolimits}
\def\Sch{\mathit{Sch}}
\def\Mor{\mathop{\mathrm{Mor}}\nolimits}
\def\Ob{\mathop{\mathrm{Ob}}\nolimits}
\def\Sh{\mathop{\mathit{Sh}}\nolimits}
\def\NL{\mathop{N\!L}\nolimits}
\def\CH{\mathop{\mathrm{CH}}\nolimits}
\def\proetale{{pro\text{-}\acute{e}tale}}
\def\etale{{\acute{e}tale}}
\def\QCoh{\mathit{QCoh}}
\def\Ker{\mathop{\mathrm{Ker}}}
\def\Im{\mathop{\mathrm{Im}}}
\def\Coker{\mathop{\mathrm{Coker}}}
\def\Coim{\mathop{\mathrm{Coim}}}

% Boxtimes
%
\DeclareMathSymbol{\boxtimes}{\mathbin}{AMSa}{"02}

%
% Macros for moduli stacks/spaces
%
\def\QCohstack{\mathcal{QC}\!\mathit{oh}}
\def\Cohstack{\mathcal{C}\!\mathit{oh}}
\def\Spacesstack{\mathcal{S}\!\mathit{paces}}
\def\Quotfunctor{\mathrm{Quot}}
\def\Hilbfunctor{\mathrm{Hilb}}
\def\Curvesstack{\mathcal{C}\!\mathit{urves}}
\def\Polarizedstack{\mathcal{P}\!\mathit{olarized}}
\def\Complexesstack{\mathcal{C}\!\mathit{omplexes}}
% \Pic is the operator that assigns to X its picard group, usage \Pic(X)
% \Picardstack_{X/B} denotes the Picard stack of X over B
% \Picardfunctor_{X/B} denotes the Picard functor of X over B
\def\Pic{\mathop{\mathrm{Pic}}\nolimits}
\def\Picardstack{\mathcal{P}\!\mathit{ic}}
\def\Picardfunctor{\mathrm{Pic}}
\def\Deformationcategory{\mathcal{D}\!\mathit{ef}}
\newenvironment{reference}{}{}
\newenvironment{slogan}{}{}
\newenvironment{history}{}{}
\usepackage{cleveref}
\crefname{originallemma}{original lemma}{original lemmas}
\Crefname{originallemma}{Original lemma}{Original lemmas}
\makeatletter
\newlabel{coherent-lemma-quasi-coherence-higher-direct-images}{{01XJ}{source}{}{sourcepointer-1}{}}
\newlabel{coherent-lemma-quasi-coherence-higher-direct-images@cref}{{[originallemma][0][]01XJ}{[1][0][]source}{}{}{}}
\makeatother
\crefname{originallemma}{original lemma}{original lemmas}
\Crefname{originallemma}{Original lemma}{Original lemmas}
\makeatletter
\newlabel{coherent-lemma-quasi-coherence-higher-direct-images-application}{{01XK}{source}{}{sourcepointer-2}{}}
\newlabel{coherent-lemma-quasi-coherence-higher-direct-images-application@cref}{{[originallemma][0][]01XK}{[1][0][]source}{}{}{}}
\makeatother
\crefname{originallemma}{original lemma}{original lemmas}
\Crefname{originallemma}{Original lemma}{Original lemmas}
\makeatletter
\newlabel{cohomology-lemma-quasi-separated-cohomology-colimit}{{01FF}{source}{}{sourcepointer-3}{}}
\newlabel{cohomology-lemma-quasi-separated-cohomology-colimit@cref}{{[originallemma][0][]01FF}{[1][0][]source}{}{}{}}
\makeatother
\crefname{originallemma}{original lemma}{original lemmas}
\Crefname{originallemma}{Original lemma}{Original lemmas}
\makeatletter
\newlabel{lemma-h1-fibre-zero}{{0E7G}{source}{}{sourcepointer-4}{}}
\newlabel{lemma-h1-fibre-zero@cref}{{[originallemma][0][]0E7G}{[1][0][]source}{}{}{}}
\makeatother
\crefname{originallemma}{original lemma}{original lemmas}
\Crefname{originallemma}{Original lemma}{Original lemmas}
\makeatletter
\newlabel{limits-lemma-higher-direct-images-zero-above-dimension-fibre}{{0E7D}{source}{}{sourcepointer-5}{}}
\newlabel{limits-lemma-higher-direct-images-zero-above-dimension-fibre@cref}{{[originallemma][0][]0E7D}{[1][0][]source}{}{}{}}
\makeatother
\title{Higher direct images of quasi-coherent sheaves vanish past degree one for proper maps with at-most-curve fibres, and the first direct image vanishes for globally generated sheaves}
\author{The Stacks Project Authors}
\date{stacks-project:more-morphisms.tex:local-snapshot-20261011; GNU FDL 1.2 or later}
\begin{document}
\maketitle
\noindent\textit{Editorial note.} Statement span is lines 20970-20986 of sources/stacks-more-morphisms/pinned-original-source.tex; proof span is lines 20988-21015, the complete author proof (26 body lines).\par\medskip
\noindent\textit{Editorial note.} Navigation stubs: 5 cited-but-not-included labels are supplied as non-load-bearing pointers with their pinned tags where the snapshot tag index has them.\par\medskip
\noindent\textit{Editorial note.} extra\_packages is empty on purpose: the retained preamble already loads xy with options and a second declaration clashes in the assembler wrapper.\par\medskip
\noindent\textit{Editorial note.} Text is retained verbatim, including the author's own typos and placeholders.\par\medskip
\noindent\textit{Editorial note.} The author's sentence contains a typo (remains is the prove (4)); retained verbatim.\par\medskip
\noindent\textit{Allowed standard prerequisites (editorial).}\begin{itemize}
\item Explicit prerequisite boundary. The dimension-bound on cohomological amplitude (Limits, lemma-higher-direct-images-zero-above-dimension-fibre), the chapter's lemma-h1-fibre-zero, quasi-coherence of higher direct images, and colimit commutation for H\textasciicircum{}1 on quasi-separated schemes.
\end{itemize}\medskip
\ifdefined\hypertarget\hypertarget{sourcepointer-1}{}\fi
\noindent\textit{Source reference.} Original lemma 01XJ: sources/stacks-more-morphisms/pinned-original-source.tex, source lines 21015--21015. Cited by the retained statement or proof but not part of the collected spans; supplied only as a navigation pointer. The argument does not depend on its content.\par\medskip
\ifdefined\hypertarget\hypertarget{sourcepointer-2}{}\fi
\noindent\textit{Source reference.} Original lemma 01XK: sources/stacks-more-morphisms/pinned-original-source.tex, source lines 21015--21015. Cited by the retained statement or proof but not part of the collected spans; supplied only as a navigation pointer. The argument does not depend on its content.\par\medskip
\ifdefined\hypertarget\hypertarget{sourcepointer-3}{}\fi
\noindent\textit{Source reference.} Original lemma 01FF: sources/stacks-more-morphisms/pinned-original-source.tex, source lines 21015--21015. Cited by the retained statement or proof but not part of the collected spans; supplied only as a navigation pointer. The argument does not depend on its content.\par\medskip
\ifdefined\hypertarget\hypertarget{sourcepointer-4}{}\fi
\noindent\textit{Source reference.} Original lemma 0E7G: sources/stacks-more-morphisms/pinned-original-source.tex, source lines 20943--20949. Cited by the retained statement or proof but not part of the collected spans; supplied only as a navigation pointer. The argument does not depend on its content.\par\medskip
\ifdefined\hypertarget\hypertarget{sourcepointer-5}{}\fi
\noindent\textit{Source reference.} Original lemma 0E7D: sources/stacks-more-morphisms/pinned-original-source.tex, source lines 21015--21015. Cited by the retained statement or proof but not part of the collected spans; supplied only as a navigation pointer. The argument does not depend on its content.\par\medskip
\begin{lemma}
\label{lemma-globally-generated-vanishing}
Let $f : X \to Y$ be a proper morphism of schemes such
that $\dim(X_y) \leq 1$ and $H^1(X_y, \mathcal{O}_{X_y}) = 0$
for all $y \in Y$. Let $\mathcal{F}$ be quasi-coherent on $X$. Then
\begin{enumerate}
\item $R^pf_*\mathcal{F} = 0$ for $p > 1$, and
\item $R^1f_*\mathcal{F} = 0$ if there is a surjection
$f^*\mathcal{G} \to \mathcal{F}$ with $\mathcal{G}$ quasi-coherent
on $Y$.
\end{enumerate}
If $Y$ is affine, then we also have
\begin{enumerate}
\item[(3)] $H^p(X, \mathcal{F}) = 0$ for $p \not \in \{0, 1\}$, and
\item[(4)] $H^1(X, \mathcal{F}) = 0$ if $\mathcal{F}$ is globally generated.
\end{enumerate}
\end{lemma}

\begin{proof}
The vanishing in (1) is Limits, Lemma
\ref{limits-lemma-higher-direct-images-zero-above-dimension-fibre}.
To prove (2) we may work locally on $Y$ and assume $Y$ is affine.
Then $R^1f_*\mathcal{F}$ is the quasi-coherent module on $Y$
associated to the module $H^1(X, \mathcal{F})$.
Here we use that $Y$ is affine, quasi-coherence of higher direct
images (Cohomology of Schemes, Lemma
\ref{coherent-lemma-quasi-coherence-higher-direct-images}), and
Cohomology of Schemes, Lemma
\ref{coherent-lemma-quasi-coherence-higher-direct-images-application}.
Since $Y$ is affine, the quasi-coherent module $\mathcal{G}$
is globally generated, and hence so is $f^*\mathcal{G}$
and $\mathcal{F}$.
In this way we see that (4) implies (2).
Part (3) follows from (1) as well as the remarks on
quasi-coherence of direct images just made. Thus
all that remains is the prove (4).
If $\mathcal{F}$ is globally generated, then there is a surjection
$\bigoplus_{i \in I} \mathcal{O}_X \to \mathcal{F}$. By part (1)
and the long exact sequence of cohomology this
induces a surjection on $H^1$. Since $H^1(X, \mathcal{O}_X) = 0$
because $R^1f_*\mathcal{O}_X = 0$ by
Lemma \ref{lemma-h1-fibre-zero}, and
since $H^1(X, -)$ commutes with direct sums
(Cohomology, Lemma \ref{cohomology-lemma-quasi-separated-cohomology-colimit})
we conclude.
\end{proof}

\end{document}
